\begin{definition}[Successor-Closed Subset]
\label{def:successor-closed-subset}
Let $(P,S,1)$ be a Peano system, and let $A \subseteq P$. The subset $A$ is
called \emph{successor-closed} exactly when
\[
\forall n \in A,\qquad S(n) \in A.
\]
\end{definition}
\end{definitionbox}

\begin{remark*}[Standard quantified statement]
For a Peano system $(P,S,1)$ and a subset $A\subseteq P$,
\[
\forall n \in A,\ S(n) \in A.
\]
\end{remark*}

\begin{remark*}[Negated quantified statement]
\[
\begin{aligned}
&\text{Let } (P,S,1) \text{ be a Peano system and let } A \subseteq P.\\
&A \text{ is not successor-closed}
\Longleftrightarrow
\exists n \in A,\ S(n) \notin A.
\end{aligned}
\]
\end{remark*}

\begin{remark*}[Failure modes]
\begin{description}
  \item[Exposition.]
  A subset fails to be successor-closed exactly when it contains some element
  of $P$ but fails to contain that element's successor.
  \item[Quantified failure.]
  \[
  \begin{aligned}
  A \text{ is not successor-closed}
  \Longleftrightarrow
  \exists n \in A,\ \underbrace{S(n) \notin A}_{\text{closure fails at } n}.
  \end{aligned}
  \]
\end{description}
\end{remark*}

\begin{remark*}[Interpretation]
A successor-closed subset is stable under one step of counting. Once an
element of $A$ is present, its successor remains in $A$ as well.
\end{remark*}

\begin{remark*}[Source comparison]
The notion of a successor-closed subset is implicit in condition (iii) of
Feferman's Definition 3.1. The present notes isolate it as a separate
definition so that induction arguments can refer to the closure condition
directly. \citep{FefermanNumberSystems1964}
\end{remark*}

\begin{dependencies}
\begin{itemize}
  \item \hyperref[def:peano-system]{Peano System}
\end{itemize}
\end{dependencies}

